\chapter{(动作行为)语法相关}

\section{行为规范}
\subsection{注释}
Python中单行注释以 \codeinline{\#} 开头
\begin{lstlisting}[style=pythonstyle]
# 这是一个注释
print("hello world!")
\end{lstlisting}

多行注释使用\codeinline{'''...'''}

\subsection{行与缩进}
python最具特色的就是使用缩进来表示代码块,
缩进的空格数是可变的，但是\textcolor{red}{\textbf{同一个代码块的语句必须包含相同的缩进空格数}}。
\begin{lstlisting}[style=pythonstyle]
if True:
    print ("True")
else:
    print ("False")
\end{lstlisting}    

\subsection{多行语句}
Python 通常是一行写完一条语句，但如果语句很长，我们可以使用反斜杠 \codeinline{\textbackslash} 来实现多行语句
\begin{lstlisting}[style=pythonstyle]
item_one = 1
item_two = 2
item_three = 3
total = item_one + \
        item_two + \
        item_three
print(total) # 输出为 6
\end{lstlisting}  

\section{运算符}
python中常用的运算符有：
\begin{table}[h]
\centering
\begin{tabular}{|l|l|}
\hline
\multicolumn{1}{|c|}{类型}  & \multicolumn{1}{c|}{包含的运算符}                                                                                                                 \\ \hline
算术运算符                   & \begin{tabular}[c]{@{}l@{}}**(指数)\\ ~(按位翻转)\\ +(一元加号)\\ -(一元减号)\\ *(乘)\\ /(除)\\ \%(取模)\\ //(取整除)\\ +(加)\\ -(减)\\ \textless\textless(左移)\\ \textgreater\textgreater(右移)\end{tabular} \\ \hline
位运算符                     & \begin{tabular}[c]{@{}l@{}}\&(按位与)\\ |(按位或)\\ \textasciicircum(按位异或)\\ \textasciitilde(按位取反)\end{tabular}                                                                 \\ \hline
赋值运算符                   & \begin{tabular}[c]{@{}l@{}}=(基础赋值)\\ +=(加后赋值)\\ -=(减后赋值)\\ *=(乘后赋值)\\ /=(除后赋值)\\ \%=(取模后赋值)\\ //=(取整除后赋值)\\ **=(指数后赋值)\\ \textless\textless=(左移后赋值)\\ \textgreater\textgreater=(右移后赋值)\end{tabular} \\ \hline
逻辑运算符                   & \begin{tabular}[c]{@{}l@{}}not(逻辑非)\\ and(逻辑与)\\ or(逻辑或)\end{tabular}                                                                                      \\ \hline
关系运算符                   & \begin{tabular}[c]{@{}l@{}}\textless(小于)\\ \textless=(小于等于)\\ \textgreater(大于)\\ \textgreater=(大于等于)\\ ==(等于)\\ != 或 \textless\textgreater(不等于)\end{tabular} \\ \hline
身份运算符                   & \begin{tabular}[c]{@{}l@{}}is(对象同一性)\\ is not(对象非同一性)\end{tabular}                                                                                 \\ \hline
成员运算符                   & \begin{tabular}[c]{@{}l@{}}in(存在于序列)\\ not in(不存在于序列)\end{tabular}                                                                                 \\ \hline
其他运算符                    & \begin{tabular}[c]{@{}l@{}}sizeof(Python 3.11+新增)\\ ?:(三元运算符)\\ .(属性访问)\\ \string[ \string](列表/字典访问)\\ ( )(元组/函数调用)\\ @(矩阵乘法)\\ ,(元组拼接)\end{tabular}                                                                 \\ \hline
\end{tabular}
\caption{Python运算符分类(基于Python 3.x语法)}
\label{tab:python_operators}
\end{table}

\subsection{算术运算符}
\begin{enumerate}[label=\arabic*)]
    \item \codeinline{+} 加法
    \item \codeinline{-} 减法
    \item \codeinline{*} 乘法
    \item \codeinline{/} 除法
    \item \codeinline{//} 整除
    \item \codeinline{\%} 取余
    \item \codeinline{**} 幂
\end{enumerate}

\begin{lstlisting}[style=pythonstyle]
x = 5
y = 2
print(x + y)
print(x - y)
print(x * y)
print(x / y)
print(x // y)
print(x % y)
print(x ** y)
\end{lstlisting}
\begin{lstlisting}[style=demostyle,numbers=none]
7
3
10
2.5
2
1
25
\end{lstlisting}

其中\codeinline{+}还可以拼接字符串, \codeinline{*}可以重复字符串
\begin{lstlisting}[style=pythonstyle]
x = 'abc'
y = 'def'
print(x + y)
print(x * 3)
\end{lstlisting}
\begin{lstlisting}[style=demostyle,numbers=none]
abcdef
abcabcabc
\end{lstlisting}

\subsection{赋值运算符}
\begin{enumerate}[label=\arabic*)]
    \item \codeinline{=} 简单赋值
    \item \codeinline{+=} 加法赋值
    \item \codeinline{-=} 减法赋值
    \item \codeinline{*=} 乘法赋值
    \item \codeinline{/=} 除法赋值
    \item \codeinline{//=} 整除赋值
    \item \codeinline{\%=} 取余赋值
    \item \codeinline{**=} 幂赋值
    \item \codeinline{|=} 或赋值
    \item \codeinline{\&=} 与赋值
    \item \codeinline{\textasciicircum{}=} 异或赋值
    \item \codeinline{>>=} 右移赋值
    \item \codeinline{<<=} 左移赋值
\end{enumerate}
\begin{lstlisting}[style=pythonstyle]
x = 5
x += 2
print("x += 2:", x)
x -= 2
print("x -= 2:", x)
x *= 2
print("x *= 2:", x)
x /= 2
print("x /= 2:", x)
x //= 2
print("x //= 2:", x)
x **= 2
print("x **= 2:", x)
x %= 2
print("x %= 2:", x)

y = 16
y |= 2
print("y |= 2:", y)
y=16
y &= 4
print("y &= 2:", y)
y=16
y >>= 2
print("y >>= 2:", y)
y <<= 2
print("y <<= 2:", y)
y ^= 2
print("y ^= 2:", y)
\end{lstlisting}
\begin{lstlisting}[style=demostyle,numbers=none]
x += 2: 7
x -= 2: 5
x *= 2: 10
x /= 2: 5.0
x //= 2: 2.0
x **= 2: 4.0
x %= 2: 0.0
y |= 2: 18
y &= 2: 0
y >>= 2: 4
y <<= 2: 16
y ^= 2: 18
\end{lstlisting}

\subsection{位运算符}
\begin{enumerate}[label=\arabic*)]
    \item \codeinline{\&} 与
    \item \codeinline{|} 或
    \item \codeinline{\textasciicircum{}} 异或
    \item \codeinline{\textasciitilde{}} 取反
    \item \codeinline{<<} 左移
    \item \codeinline{>>} 右移
\end{enumerate}
\begin{lstlisting}[style=pythonstyle]
x = 16
y = 2
print("x & y:", x & y)
print("x | y:", x | y)
print("x ^ y:", x ^ y)
print("~x:", ~x)
print("x << 2:", x << 2)
print("x >> 2:", x >> 2)
\end{lstlisting}    
\begin{lstlisting}[style=demostyle,numbers=none]  
x & y: 0
x | y: 18
x ^ y: 18
~x: -17
x << 2: 64
x >> 2: 4    
\end{lstlisting}    

\subsection{比较运算符}
\begin{enumerate}[label=\arabic*)]
    \item \codeinline{==} 两个变量的值是否相等
    \item \codeinline{!=} 不等于
    \item \codeinline{>} 大于
    \item \codeinline{<} 小于
    \item \codeinline{>=} 大于等于
    \item \codeinline{<=} 小于等于
    \item \codeinline{is} 判断两个变量所引用的对象是否相同
    \item \codeinline{is not} 判断两个变量所引用的对象是否不相同
\end{enumerate}

\begin{lstlisting}[style=pythonstyle]
import time #引入 time 模块

t1 = time.gmtime() # gmtime()用来获取当前时间(秒)
t2 = time.gmtime()

print(t1 == t2) #输出 True
print(t1 is t2) #输出 False
\end{lstlisting}

\subsection{逻辑运算符}
\begin{enumerate}[label=\arabic*)]
    \item \codeinline{and} 与
    \item \codeinline{or} 或
    \item \codeinline{not} 非
\end{enumerate}

\begin{lstlisting}[style=pythonstyle]
x = 5
y = 2
print(x > y and x < y)
print(x > y or x < y)
print(not x > y)
\end{lstlisting} 
\begin{lstlisting}[style=demostyle,numbers=none]  
False
True
False
\end{lstlisting}


\section{逻辑}
\subsection{条件}
\begin{lstlisting}[style=pythonstyle]
x = 'abc'
if x == 'abc':
    print('x is abc')
elif x == 'def':
    print('x is def')
else:
    print('x is not abc or def')
\end{lstlisting}

\subsection{三元运算符}
\begin{lstlisting}[style=demostyle,numbers=none]  
表达式1 if 条件 else 表达式2
\end{lstlisting}

\begin{lstlisting}[style=pythonstyle]
x = 'abc'
print('x is abc') if x == 'abc' else print('x is not abc')
\end{lstlisting}
\begin{lstlisting}[style=demostyle,numbers=none] 
x is abc
\end{lstlisting}

\begin{lstlisting}[style=pythonstyle]
a = 10
b = 15
max = a if a>b else b
print("max=", max)
\end{lstlisting}
\begin{lstlisting}[style=demostyle,numbers=none] 
max= 15
\end{lstlisting}



\subsection{循环}
\subsection{while}
\begin{lstlisting}[style=pythonstyle,numbers=none]
i = 0
while i < 5:
    print(i)
    i += 1
\end{lstlisting}

\subsection{for}
Python中的\codeinline{for}和C语言中的\codeinline{for}有所不同, Python的\codeinline{for}是用来遍历序列的。

\begin{lstlisting}[style=demostyle,numbers=none]
for i in range(5):
    print(i)
\end{lstlisting}

